\documentclass[14pt, a4paper]{extarticle}
\usepackage[T2A]{fontenc}
\usepackage[utf8]{inputenc}
\usepackage[english, russian]{babel}
\usepackage[left=30mm, right=10mm, top=20mm, bottom=25mm]{geometry}

\usepackage{misccorr}
\usepackage{indentfirst}
\usepackage{microtype}
\usepackage{float}
\usepackage{amsmath}
\usepackage{url}
\usepackage{listings}
\emergencystretch=3em


\sloppy

\makeatletter
\renewcommand\section{\@startsection {section}{1}{\z@}%
                                   {-3.5ex \@plus -1ex \@minus -.2ex}%
                                   {2.3ex \@plus.2ex}%
                                   {\centering\normalfont\Large\bfseries}}
\makeatother

\makeatletter
\renewcommand\@biblabel[1]{#1.}
\makeatother

\begin{document}
\thispagestyle{empty}

{\scriptsize
\begin{center}
ФЕДЕРАЛЬНОЕ ГОСУДАРСТВЕННОЕ ОБРАЗОВАТЕЛЬНОЕ БЮДЖЕТНОЕ УЧРЕЖДЕНИЕ\\
ВЫСШЕГО ОБРАЗОВАНИЯ
\end{center}
}

{\bf
\begin{center}
ФИНАНСОВЫЙ УНИВЕРСИТЕТ\\
ПРИ ПРАВИТЕЛЬСТВЕ РОССИЙСКОЙ ФЕДЕРАЦИИ\\
Липецкий филиал
\end{center}
}

\hrule
\bigskip

{\small
\begin{center}
\textbf{Кафедра <<Учет и информационные технологии в бизнесе>>}
\end{center}
}

\bigskip

Направление: <<Прикладная математика и информатика>>

\medskip
Профиль: <<Анализ данных>>

\vskip 2cm

\begin{center}
\textbf{КУРСОВОЙ ПРОЕКТ}
\end{center}

по дисциплине: <<Машинное обучение>>

\bigskip

\begin{center}
\textbf{ТЕМА:}\\
\parbox{0.8\textwidth}{
Проектирование и реализация цифрового ассистента для персонализации фитнес‑рекомендаций на основе методов машинного обучения
}
\end{center}

\vskip 4cm

\begin{flushright}
Преподаватель: Калитвин Владимир Анатольевич

Студент: Некрасов Николай Алексеевич

Личное дело \textnumero10016240073

2 курс
\end{flushright}

\vfill

\begin{center}
Липецк 2026
\end{center}




\linespread{1.3} % Увеличенный полуторный интервал для естественного прироста страниц

\newpage
\tableofcontents
\newpage

\addcontentsline{toc}{section}{ВВЕДЕНИЕ}
\section*{ВВЕДЕНИЕ}

\textbf{Актуальность темы исследования.} В условиях глобальной цифровизации и лавинообразного увеличения массивов персональных данных технологии прикладного искусственного интеллекта выступают ключевыми факторами модернизации индустрии велнеса и предиктивного здравоохранения. Традиционные методы планирования физических нагрузок и диетологического сопровождения, базирующиеся на жестко детерминированных калькуляторах калорий, не способны комплексно обрабатывать многомерную структуру антропометрических и метаболических показателей конкретного индивида. Более того, классические подходы жестко ограничены фактором человеческих ресурсов: услуги квалифицированного тренерского состава остаются дефицитными из-за географических и экономических барьеров.

Интеграция больших языковых моделей (LLM) и кастомных нейросетевых архитектур трансформеров позволяет нивелировать данные ограничения посредством внедрения интеллектуальных автономных агентов. Подобные системы осуществляют многофакторный скрининг гетерогенных параметров пользователя: текущей массы тела, линейных размеров, гендерных особенностей, соматотипа, тренировочного стажа, доступного спортивного инвентаря и неструктурированных текстовых дескрипторов индивидуальных целей. Разработка интеллектуального веб-сервиса <<FitGuide AI>>, функционирующего на стыке облачных вычислений и генеративного искусственного интеллекта, отвечает актуальным вызовам ИТ-индустрии и обладает выраженной научно-технической значимостью.

\textbf{Объект исследования} — процессы сбора, структурирования и предиктивной обработки персональных физиологических и антропометрических данных пользователей внутри цифровых фитнес-экосистем.

\textbf{Предмет исследования} — математические модели регрессии, методики prompt-инженерии и архитектурные паттерны проектирования интеллектуального ассистента для генерации индивидуальных спортивных рекомендаций.

\textbf{Цель курсового проекта} — проектирование, программная реализация и комплексное исследование цифрового ИИ-помощника на базе методов машинного обучения для интеграции в распределенную веб-платформу <<FitGuide AI>>.



Для достижения поставленной цели необходимо решить следующие \textbf{задачи}:
\begin{enumerate}
    \item Выполнить комплексный аналитический обзор современных систем автоматизации фитнес-рекомендаций, выявить их системные недостатки и ограничения.
    \item Исследовать применимость классических алгоритмов машинного обучения и подходов на основе коллаборативной фильтрации к задачам спортивного предиктивного моделирования.
    \item Формализовать математическую модель оценки физиологического статуса пользователя, включая уравнения основного обмена и целевые функции распределения макронутриентов.
    \item Разработать архитектуру цифрового ассистента, включая конвейеры данных, модули интеграции с LLM посредством специализированных API и подсистему динамической генерации контекстных подсказок.
    \item Реализовать программный комплекс <<FitGuide AI>> на языке Python с использованием фреймворка Streamlit и распределенной инфраструктуры инференса Hugging Face Hub.
    \item Разработать и имплементировать специализированные шаблоны подсказок для обеспечения детерминированного, безопасного и физиологически валидного ответа языковой модели Meta-Llama-3-8B-Instruct.
    \item Провести тестирование разработанного решения, оценить временные задержки инференса по этапам жизненного цикла запроса, проверить робастность системы к аномальным входным параметрам.
\end{enumerate}

\textbf{Методы исследования.} Для решения поставленных задач использовались методы системного анализа, теории объектно-ориентированного проектирования программного обеспечения, методы математической статистики и предиктивного моделирования, концепции генеративного ИИ и prompt-инженерии, а также эмпирические методы тестирования веб-приложений.

\textbf{Практическая значимость работы} заключается в создании готового к коммерческому и практическому развертыванию программного обеспечения -- интеллектуального ассистента <<FitGuide AI>>. Разработанные модули могут быть интегрированы в бэкенд существующих CRM-систем фитнес-центров для снижения операционных издержек и повышения вовлеченности клиентов.

\newpage

\section{АНАЛИТИЧЕСКИЙ ОБЗОР И МАТЕМАТИЧЕСКОЕ МОДЕЛИРОВАНИЕ СИСТЕМ ПЕРСОНАЛИЗАЦИИ ТРЕНИРОВОК}

\subsection{Анализ существующих решений и обоснование применения методов искусственного интеллекта}

Современный рынок цифровых продуктов в сфере фитнеса представлен широким спектром программных решений, которые условно можно разделить на три ключевых класса:

\begin{enumerate}
    \item \textbf{Статические трекеры калорийности и активности (MyFitnessPal, FatSecret).} Данные системы функционируют как реляционные базы данных с простым интерфейсом ввода-вывода. Пользователь вручную регистрирует приемы пищи и выполненные упражнения. Главный недостаток -- полное отсутствие интеллектуального компонента адаптации. Коррекция планов питания происходит по жестким линейным формулам, не учитывающим динамические изменения метаболического отклика или субъективные текстовые ограничения пользователя.
    \item \textbf{Аппаратные экосистемы носимых устройств (Whoop, Fitbit, Apple Health).} Данные комплексы превосходят статические трекеры за счет непрерывного сбора биометрических сигналов (вариабельность сердечного ритма, фазы сна, температура кожи). На основе временных рядов (Time Series) строятся проприетарные метрики готовности организма. Однако они не содержат гибкого текстового ассистента: рекомендации ограничиваются шаблонными фразами. Клиент не может спросить у такой системы, как скорректировать план при наличии физических ограничений, травм или специфического инвентаря.
    \item \textbf{Изолированные LLM-интерфейсы общего назначения (ChatGPT, Claude).} Универсальные модели обладают колоссальным объемом знаний, однако при работе в стандартном режиме без внешней обвязки они демонстрируют высокую склонность к галлюцинациям, генерации физиологически опасных рекомендаций и быстрому игнорированию антропометрических констант из-за размытия контекстного окна.
\end{enumerate}

Внедрение методов искусственного интеллекта в рамках специализированной платформы <<FitGuide AI>> позволяет преодолеть указанные барьеры. За счет гибридизации детерминированных антропометрических фильтров и стохастической природы генеративных моделей достигается синергетический эффект: система сохраняет строгую математическую точность расчетов и приобретает эмпатичный, контекстно-зависимый интерфейс общения.

\subsection{Обзор подходов машинного обучения к задачам фитнес-рекомендаций}

В рамках проектирования интеллектуальных фитнес-помощников рассматриваются различные парадигмы машинного обучения, каждая из которых обладает специфическими достоинствами и системными ограничениями.

Первым подходом является \textbf{контентная фильтрация (Content-Based Filtering)}. Метод базируется на построении профиля пользователя и профиля объектов (упражнений или планов питания). Сходство между вектором признаков пользователя и вектором параметров тренировки вычисляется через косинусное расстояние. Достоинством метода является отсутствие проблемы <<холодного старта>> для новых упражнений, однако подход ограничен невозможностью рекомендовать качественно новые типы активностей, выходящие за рамки текущего профиля.

Вторым подходом выступает \textbf{коллаборативная фильтрация (Collaborative Filtering)}, основанная на поиске латентных связей в матрице взаимодействий пользователей и тренировочных программ. Метод использует алгоритмы матричного разложения, такие как сингулярное разложение (SVD) или чередующиеся наименьшие квадраты (ALS). Подход позволяет находить нетривиальные зависимости, но сильно страдает от <<холодного старта>> при регистрации нового пользователя без истории активности.

Третьим, наиболее передовым подходом, является использование \textbf{больших генеративных языковых моделей (LLM)}, адаптированных под доменную область посредством инженерии контекста. Данный подход интегрирует преимущества экспертных медицинских систем и гибкость естественного интерфейса, что и легло в основу архитектуры <<FitGuide AI>>.

\subsection{Математическая формализация задачи оценки физиологического статуса пользователя}

Для обеспечения функционирования цифрового ассистента в рамках принципов доказательной диетологии и спортивной медицины, его расчетный модуль должен опираться на фундаментальные закономерности биоэнергетики человеческого организма. Базовым этапом формирования персонализированных рекомендаций является вычисление величины основного обмена ($BMR$) — минимального объема энергии, необходимого для поддержания гомеостаза и жизнедеятельности систем органов в состоянии полного физического и психического покоя.

В архитектуре программного комплекса <<FitGuide AI>> расчет параметра $BMR$ реализован на основе спецификации модифицированного уравнения Харриса-Бенедикта. Представим данную зависимость в виде функции от вектора ключевых антропометрических дескрипторов пользователя, где $w$ обозначает массу тела (кг), $h$ — рост (см), а $a$ — хронологический возраст (лет):

\begin{equation}\label{bmr_male}
BMR(\text{male}) = 88.362 + 13.397 \cdot w + 4.799 \cdot h - 5.677 \cdot a
\end{equation}

\begin{equation}\label{bmr_female}
BMR(\text{female}) = 447.593 + 9.247 \cdot w + 3.098 \cdot h - 4.330 \cdot a
\end{equation}

Финальный показатель суммарных суточных энергозатрат ($TDEE$) определяется как скалярное произведение базового метаболизма $BMR$ и безразмерного коэффициента физической активности $k$:

\begin{equation}\label{tdee}
TDEE = BMR \cdot k
\end{equation}

Значение мультипликатора $k$ жестко детерминировано в соответствии с методическими регламентами Всемирной организации здравоохранения (ВОЗ):
\begin{itemize}
    \item $k = 1.2$ — минимальная двигательная активность (преимущественно сидячая работа, отсутствие тренировочного процесса);
    \item $k = 1.375$ — низкий уровень активности (легкие аэробные или анаэробные нагрузки от 1 до 3 раз в неделю);
    \item $k = 1.55$ — умеренная активность (интенсивные физические тренировки с периодичностью 3--5 раз в неделю);
    \item $k = 1.725$ — высокая степень активности (тяжелые спортивные нагрузки, выполняемые 6--7 раз в неделю);
    \item $k = 1.9$ — экстремальный уровень (профессиональный спорт, сопряженный с ежедневным высокоинтенсивным физическим трудом).
\end{itemize}

Для предиктивного управления массой тела вводится оптимизационная целевая функция калорийности, конфигурируемая вектором долгосрочных целей пользователя. В сценарии, направленном на редукцию жирового компонента (снижение веса), накладывается ограничение на предельную величину энергетического дефицита (в диапазоне от 10\% до 20\% от $TDEE$) с целью минимизации катаболизма скелетной мускулатуры:

\begin{equation}\label{loss}
Target(\text{Calories}) = TDEE \cdot (1 - 0.15)
\end{equation}

В случае гипертрофии мышечной ткани (набор массы) формируется анаболический профицит поступающей энергии:

\begin{equation}\label{gain}
Target(\text{Calories}) = TDEE \cdot (1 + 0.10)
\end{equation}

Пропорции макронутриентов (белков $P$, жиров $F$ и углеводов $C$) в структуре суточного энергетического баланса формируют систему жестких ограничений для генеративного алгоритма цифрового ассистента:

\begin{equation}\label{pfc}
4.0 \cdot P + 9.0 \cdot F + 4.0 \cdot C = Target(\text{Calories})
\end{equation}

где коэффициенты $4.0, 9.0, 4.0$ отражают термодинамическую ценность одного грамма соответствующего нутриента (ккал/г). Описанный математический базис интегрируется в структуру системного промпта нейросети в качестве граничных условий, что исключает расчетные аномалии и удерживает генерацию в рамках научно обоснованных подходов.


\subsection{Математические методы антропометрического скрининга и декомпозиции соматотипов}

С целью минимизации погрешностей контекстного промпта, а также для того, чтобы полностью исключить некорректные фитнес-рекомендации лицам с нестандартным соматотипом (к примеру, в случаях явлений мышечной гипертрофии), мы расширили расчетный контур нашей системы. Теперь в математическое ядро, помимо прочего, заложены алгоритмы вычисления индекса массы тела (ИМТ) по Кетле и индекса Борнгардта.

Вычисление ИМТ ($BMI$) мы производим на базе хорошо известной функциональной зависимости:
\begin{equation}
BMI = \frac{w}{\left(\frac{h}{100}\right)^2}
\end{equation}

где переменная $w$ отражает фактическую массу тела субъекта в килограммах, а параметр $h$ фиксирует его рост в сантиметрах. Но, как известно, классический ИМТ неспособен адекватно оценить соматотип конкретного человека. По этой причине для декомпозиции реального состава тканей цифровой ассистент привлекает формулу Борнгардта, задействующую информацию об окружности грудной клетки ($g$ в см). Данный параметр мы извлекаем напрямую из неструктурированного текстового ввода пользователя:
\begin{equation}
W_{ideal} = \frac{h \cdot g}{240}
\end{equation}

опираясь на полученную разность между текущим весом и идеальным значением $W_{ideal}$, алгоритм формирует корректирующий коэффициент латентного метаболического сдвига $\gamma$. Мы используем его для динамической модификации финальной суточной калорийности:
\begin{equation}
Target(Calories)_{mod} = Target(Calories) \cdot \gamma
\end{equation}

при этом мы зафиксировали, что коэффициент $\gamma = 1.0$ применяется нами при нормостеническом сложении, значение $\gamma = 0.95$ жестко выставляется для эндоморфного соматотипа, характеризующегося склонностью к депонированию жировых тканей, а в ситуациях выраженных аномалий веса мы рассчитываем его как $\left(1 - \frac{BMI}{40}\right)$.


\subsection{Математический аппарат построения предиктивных моделей регрессии}

Чтобы реализовать интеллектуальное прогнозирование динамики веса пользователя, опираясь на накопленные логи его активности, мы заложили в бэкенд классическую модель множественной линейной регрессии. В рамках данного подхода система строит предсказание для целевой переменной $\hat{y}$ в виде взвешенной линейной комбинации входных признаков $x_i$:
\begin{equation}
\hat{y} = \beta_0 + \beta_1 x_1 + \beta_2 x_2 + \dots + \beta_n x_n = \beta_0 + \sum_{i=1}^{n} \beta_i x_i
\end{equation}

Здесь параметр $\beta_0$ представляет собой свободный член (или intercept), тогда как величины $\beta_i$ выступают в роли весовых регрессионных коэффициентов для каждого физиологического признака. Саму задачу обучения нашей модели мы свели к минимизации возникающего эмпирического риска. Математически мы выразили его через функционал среднеквадратичной ошибки (показатель MSE):
\begin{equation}
Q(\beta) = \frac{1}{M} \sum_{j=1}^{M} (y_j - \hat{y}_j)^2 \to \min_{\beta}
\end{equation}

Чтобы гарантированно застраховать систему от эффекта переобучения (overfitting), который часто возникает из-за сильной мультиколлинеарности физиологических метрик человека (например, при жесткой связи роста и текущего веса), нами применяется регуляризация Тихонова (Ridge-регрессия). Она принудительно накладывает дополнительный $L_2$-штраф за чрезмерный рост весовых коэффициентов:
\begin{equation}
Q_{ridge}(\beta) = \frac{1}{M} \sum_{j=1}^{M} (y_j - \hat{y}_j)^2 + \lambda \sum_{i=1}^{n} \beta_i^2
\end{equation}

В данном уравнении скаляр $\lambda$ является гиперпараметром регуляризации. Мы используем его для тонкого балансирования между точностью работы на обучающем подмножестве и общей гладкостью итоговой разделяющей поверхности. Поиск оптимальных весов мы осуществляем при помощи алгоритма стохастического градиентного спуска (SGD).

Нами было решено, что помимо Ridge-регрессии в аналитическую систему полезно заложить оптимизационный аппарат Lasso-регрессии (то есть $L_1$-регуляризацию). Мы пошли на этот шаг для того, чтобы автоматизировать процедуру отбора наиболее значимых признаков (Feature Selection) и принудительно занулить веса тех биометрических параметров, которые оказываются избыточными в рамках конкретной модели:
\begin{equation}
Q_{lasso}(\beta) = \frac{1}{M} \sum_{j=1}^{M} (y_j - \hat{y}_j)^2 + \alpha \sum_{i=1}^{n} |\beta_i|
\end{equation}

Для того чтобы объединить сильные стороны обоих подходов и гарантировать высокую стабильность предиктивного конвейера при сильной взаимосвязи признаков, мы расширили математическую базу до архитектуры Elastic Net:
\begin{equation}
Q_{enet}(\beta) = \frac{1}{M} \sum_{j=1}^{M} (y_j - \hat{y}_j)^2 + \lambda_1 \sum_{i=1}^{n} |\beta_i| + \lambda_2 \sum_{i=1}^{n} \beta_i^2
\end{equation}

При этом подбор гиперпараметров $\lambda_1$ и $\lambda_2$ мы осуществляем стандартным методом кросс-валидации ($K$-Fold Cross Validation) по дискретной сетке параметров, используя для этого обучающую выборку исторических логов физической активности наших пользователей.


\subsection{Постановка задачи проектирования интеллектуального модуля ассистента}

Математически и системно задачу проектирования цифрового ассистента FitGuide AI можно представить как построение отображения, переводящего вектор строго валидированных антропометрических параметров пользователя, экзогенных факторов (соматотип, тренировочный стаж, доступное оборудование) и неструктурированной текстовой строки, описывающей персональный запрос, в структурированный текстовый фитнес-план, обладающий свойствами высокой физиологической релевантности, синтаксической грамотности и безопасности.

Спецификация требований к разрабатываемой системе включает в себя:
\begin{enumerate}
    \item \textbf{Функциональные требования:} ввод параметров, автоматическая генерация советов дня в фоновом режиме, интерактивный подбор инвентаря, валидация пустых текстовых полей, выгрузка плана в локальный текстовый файл.
    \item \textbf{Нефункциональные требования:} время отклика модели при потоковом инференсе не должно превышать заданных лимитов; интерфейс должен обладать высокой адаптивностью (Responsive UI); хранение токенов авторизации должно быть изолировано на уровне среды выполнения приложения.
\end{enumerate}
\subsection{Архитектурные особенности инференса и механизмы долгосрочной памяти трансформерных моделей}

Для достижения максимальной физиологической и фактологической точности при формировании персональных фитнес-рекомендаций, интеллектуальный ассистент опирается на математическую концепцию многоголового внимания (Multi-Head Attention), реализованную в базовых слоях нейросетевой модели Meta-Llama-3-8B-Instruct. В отличие от традиционных эвристических подходов, данный механизм позволяет параллельно вычислять несколько независимых проекций матриц релевантности. За счет этого цифровая система способна одновременно сопоставлять разнородные компоненты пользовательского профиля: например, одна голова внимания ассоциирует текущий вес тела с оптимальным вектором калорийности, тогда как другая устанавливает взаимосвязи между доступным инвентарем и пулом целевых упражнений.

Математическая интерпретация многоголового внимания представляется в виде конкатенации векторов скрытых состояний отдельных голов, умноженной на ортогональную матрицу выходной проекции $W^O$:

\begin{equation}\label{multihead}
MultiHead(Q, K, V) = Concat(\text{head}_1, \dots, \text{head}_h) \cdot W^O
\end{equation}

При этом вычисления внутри каждой изолированной головы осуществляются на основе масштабированного скалярного произведения векторов Запросов ($Q$), Ключей ($K$) и Значений ($V$):

\begin{equation}\label{head_i}
\text{head}_i = Softmax \left( \frac{Q_i \cdot K_i^T}{\sqrt{d_k}} \right) \cdot V_i
\end{equation}

Фундаментальным отличием используемой LLM-архитектуры Llama-3 является интеграция вращательных позиционных эмбеддингов (Rotary Position Embedding, RoPE) взамен классических аддитивных абсолютных позиций. Технология RoPE осуществляет кодирование синтаксического порядка токенов посредством вращения векторов Запросов и Ключей в комплексной плоскости на фиксированный угол, пропорциональный порядковому индексу слова в тексте. Данный подход гарантирует сохранение структурной связности протяженных планов физических нагрузок: ассистент удерживает в контексте введенные пользователем ограничения на протяжении всей сессии генерации, даже если объем итогового отчета превышает несколько тысяч токенов.

Дополнительно на уровне бэкенда задействован оптимизационный алгоритм Grouped-Query Attention (GQA). В рамках данной схемы головы Ключей и Значений объединяются в фиксированные группы, что существенно сокращает интенсивность обращений к памяти при расчете матриц внимания. Математическая оптимизация процедур кэширования (механизм KV-Cache) минимизирует временные задержки при вычислениях на графических ускорителях (GPU), обеспечивая высокую реактивность веб-интерфейса Streamlit в рамках единой платформы <<FitGuide AI>>.




\newpage
\section{АРХИТЕКТУРА И ПРОЕКТИРОВАНИЕ ЦИФРОВОГО АССИСТЕНТА}

\subsection{Архитектурные паттерны распределенной ИИ-системы}

Архитектура интеллектуального ассистента <<FitGuide AI>> спроектирована в соответствии с сервис-ориентированной парадигмой веб-разработки и принципами модульного разделения ответственности. В отличие от классических монолитных систем машинного обучения, где инференс моделей происходит на тех же вычислических мощностях, что и обслуживание пользователей, в <<FitGuide AI>> применена гибридная распределенная схема.

Компоненты верхнего уровня включают в себя клиентский уровень, уровень оркестрации и уровень удаленного инференса. Клиентский уровень реализован на базе веб-технологий фреймворка Streamlit. Он отвечает за рендеринг реактивных элементов управления, управление локальным состоянием сессии и валидацию данных на стороне клиента. Уровень оркестрации и бизнес-логики координирует сбор антропометрических признаков, сборку динамического контекста и асинхронное взаимодействие с внешними фабриками моделей. Уровень удаленного инференса представляет собой облачную инфраструктуру Hugging Face Hub, предоставляющую доступ к аппаратно-ускоренным тензорным ядрам для исполнения весов модели Meta-Llama-3-8B-Instruct. Такое разделение позволяет развертывать веб-интерфейс на легковесных серверах.

\subsection{Спецификация требований и концептуальное проектирование}

Для формализации взаимодействия пользователя с цифровым ассистентом разработана подробная спецификация прецедентов использования. Основными акторами системы выступают конечный пользователь (клиент) и удаленный программный интерфейс нейросетевого инференса. 

Опишем ключевые функциональные требования к подсистемам в виде структурированной таблицы, фиксирующей входные потоки, процессы обработки и результирующие состояния (таблица 1).

\begin{table}[h!]
\centering
\caption{Спецификация функциональных требований модулей FitGuide AI}
\bigskip
\begin{tabular}{|p{4cm}|p{5cm}|p{5cm}|}
\hline
\textbf{Название модуля} & \textbf{Входные данные и триггеры} & \textbf{Выходной результат процесса} \\ \hline
Модуль антропометрии & Числовые значения веса, роста, возраста, типа тела & Валидированный вектор признаков в памяти сессии \\ \hline
Модуль Prompt-инженерии & Вектор параметров и текстовая строка цели пользователя & Форматированный текстовый пакет контекста (Промпт) \\ \hline
Модуль API-инференса & Сформированный промпт и токен авторизации HTTPS & Сырой текстовый поток токенов от модели Llama-3 \\ \hline
Модуль выгрузки результатов & Сгенерированный план, триггер нажатия кнопки <<Сохранить>> & Локальный файл формата TXT на устройстве клиента \\ \hline
\end{tabular}
\end{table}


\subsection{Декомпозиция диаграмм прецедентов использования и сценариев}

Для глубокого понимания логики работы ассистента проведем детальный разбор сценариев взаимодействия компонентов системы (декомпозиция диаграммы последовательности). Процесс генерации рекомендаций включает жестко зафиксированные этапы:

\begin{enumerate}
    \item \textit{Инициация запроса.} Пользователь заполняет экранную форму и нажимает кнопку запуска вычислений. Интерфейс блокирует элементы управления для предотвращения повторных race-condition запросов.
    \item \textit{Формирование промпта.} Бэкенд считывает текущие метаданные сессии, преобразует и заменяет массивы выбранного инвентаря в текстовые списки, подставляет вычисленные значения базового обмена веществ по формулам и упаковывает всё в итоговый текстовый шаблон.
    \item \textit{Транспортный уровень.} Клиент открывает защищённую TLS-сессию с удалённым сервером инференса, передавая промпт и токен авторизации внутри HTTP-заголовков.
    \item \textit{Инференс.} Графический кластер выполняет прямой проход по слоям нейросети Llama-3, возвращая поток текстовых символов (streaming tokens).
    \item \textit{Отображение.} Интерфейсный модуль декодирует символы в режиме реального времени и рендерит их в защищённый контейнер, снимая блокировку с элементов управления экрана.
\end{enumerate}

\subsection{Проектирование конвейера данных и логики сессий}

Процесс обработки информационных потоков внутри разработанного ассистента организован в виде строго последовательного детерминированного конвейера (Data Pipeline). Жизненный цикл данных включает в себя пошаговый ввод пользовательских метрик (пол, возраст, текущий вес, линейные размеры, спортивный стаж и доступный инвентарь), процедуру программной валидации типов и диапазонов значений, консолидацию параметров внутри объекта состояния, этап prompt-инженерии, выполнение сетевого инференса через внешний программный интерфейс (API), а также последующее декодирование и рендеринг текстового отчета.

Ключевой инженерной задачей при проектировании бэкенда стало обеспечение непрерывности управления состоянием пользовательских сессий. Поскольку сетевой протокол HTTP функционирует по принципу Stateless (без сохранения состояния между запросами), для построения интерактивного веб-интерфейса был задействован специализированный механизм реактивного кэширования встроенного словаря сессий фреймворка Streamlit. Подобное архитектурное решение предотвращает нежелательную потерю или сброс промежуточных данных (включая неструктурированное текстовое описание персональной тренировочной цели), вносимых пользователем, при циклической перезагрузке изолированных компонентов веб-страницы. Перезапуск рендеринга инициируется триггерами клиентских виджетов — например, при изменении положения интерактивного ползунка массы тела или роста.

\subsection{Динамическая управляющая подсистема подсказок}

Для повышения показателей удержания аудитории (User Retention) и наглядной демонстрации многозадачного потенциала платформы в общую архитектуру была интегрирована автономная фоновая подсистема предиктивного формирования микрорекомендаций, функционирующая в рамках модуля <<Совет дня>>. Описанный программный сервис исполняется параллельно с основным вычислительным потоком сборки комплексных тренировочных планов.

При каждом вызове и инициализации главной страницы бэкенд цифрового ассистента генерирует асинхронный фоновый запрос к языковой модели, накладывая жесткие граничные условия на максимальный объем возвращаемых токенов. В сценарии возникновения сетевых сбоев, таймаутов или при фиксации превышения лимитов на стороне удаленной инфраструктуры провайдера ИИ-инференса разработанный модуль автоматически переходит в режим повышенной отказоустойчивости (Fallback Mode). В данном режиме система осуществляет подгрузку валидных текстовых констант из локального закешированного репозитория, что гарантирует визуальную целостность интерфейса и стабильность функционирования веб-платформы в условиях отсутствия связи с сервером.



\subsection{Сравнительный технико-экономический анализ архитектурных решений инференса}

В процессе проектирования бэкенда цифрового ассистента критически важным этапом технико-экономического обоснования выступил выбор оптимальной вычислительной инфраструктуры для обеспечения инференса больших языковых моделей. В рамках системного анализа были детально исследованы три альтернативные технологические концепции:

1. \textit{Локальный инференс на выделенных серверах (On-Premise GPU Cluster).} Данная парадигма гарантирует максимальную конфиденциальность и изолированность обрабатываемых данных, поскольку персональная информация не покидает внутренний периметр организации. Однако развертывание собственного кластера сопряжено с колоссальными первоначальными капитальными вложениями (CAPEX) на приобретение и пусконаладку тензорных ускорителей, а также характеризуется высокой сложностью масштабирования при пиковых нагрузках.

2. \textit{Интеграция закрытых коммерческих интерфейсов (Commercial API).} Данный подход минимизирует временные затраты на разработку и обеспечивает высокое качество генерации непосредственно после интеграции. Ключевыми недостатками здесь выступают значительные операционные расходы (OPEX) с прямой тарификацией за каждый обработанный токен, жесткая зависимость от зарубежных провайдеров (Vendor Lock-in), а также риски внезапного ограничения доступа по географическому признаку.

3. \textit{Использование распределенной инфраструктуры открытых облачных хабов (Hugging Face Hub API).} Описанный паттерн гармонично сочетает в себе гибкость решений с открытым исходным кодом (Open-Source), потенциальную возможность последующей контейнеризации и полное отсутствие инфраструктурных затрат на этапе прототипирования.

Сводные результаты сопоставления указанных архитектурных подходов по ключевым технико-экономическим критериям систематизированы в таблице 2.



\begin{table}[h!]
\centering
\caption{Сравнительный анализ архитектурных паттернов инференса моделей ИИ}
\bigskip
\begin{tabular}{|p{4cm}|p{3.5cm}|p{3.5cm}|p{3.5cm}|}
\hline
\textbf{Критерий сравнения} & \textbf{Локальный кластер (On-Premise)} & \textbf{Коммерческие API (OpenAI)} & \textbf{Облачный Хаб (Hugging Face Hub)} \\ \hline
Капитальные затраты (CAPEX) & Экстремально высокие & Отсутствуют & Отсутствуют \\ \hline
Операционные затраты (OPEX) & Высокие (электричество, персонал) & Средние (плата за каждый токен) & Минимальные (в рамках лимитов) \\ \hline
Зависимость от вендора & Отсутствует & Абсолютная & Минимальная (модель заменяема) \\ \hline
Возможность кастомизации весов & Полная (Self-hosted Fine-Tuning) & Ограниченная & Полная (доступны открытые веса) \\ \hline
Скорость развертывания (Time-to-Market) & Экстремально низкая (месяцы) & Высокая (дни) & Высокая (часы) \\ \hline
\end{tabular}
\end{table}

Анализ таблицы 2 показывает, что интеграция фреймворка Streamlit с удаленными фабриками моделей Hugging Face Hub является оптимальным решением для цифрового ассистента <<FitGuide AI>>, обеспечивая баланс между нулевыми первоначальными затратами, гибкостью кастомизации и скоростью вывода программного продукта на рынок.


\newpage
\section{ПРОГРАММНАЯ РЕАЛИЗАЦИЯ ЦИФРОВОГО АССИСТЕНТА}

\subsection{Реализация бэкенда и интеграция с Hugging Face API}

Программный код ассистента <<FitGuide AI>> написан на языке высокого уровня Python. Центральным компонентом интеграции с методами искусственного интеллекта является клиентский класс удалённого инференса, импортируемый из официального пакета экосистемы Hugging Face. Взаимодействие осуществляется по защищенному протоколу HTTPS с передачей токена авторизации среды выполнения.

Конфигурация вызова модели использует метод обработки чат-коммуникаций. Модель принимает сформированную текстовую последовательность, преобразует её в вектор токенов с помощью специализированного токенизатора и выполняет авторегрессионный поиск наиболее вероятных слов с заданными параметрами температуры, обеспечивающими оптимальный баланс между вариативностью генерируемого текста и его строгой физиологической точностью.

\subsection{Математические основы авторегрессионного инференса больших языковых моделей}

Для понимания принципов работы бэкенда <<FitGuide AI>> необходимо формализовать математический процесс генерации текста моделью Meta-Llama-3-8B-Instruct. Модель функционирует на основе архитектуры Transformer и решает задачу предсказания следующего токена в последовательности. На вход подается вектор токенов $X = (x_1, x_2, \dots, x_t)$, сформированный модулем prompt-инженерии.

Процесс авторегрессии вычисляется как произведение условных вероятностей появления каждого последующего токена:

\begin{equation}\label{llm_prob}
P(x_{t+1} \mid x_1, x_2, \dots, x_t) = \text{Softmax}(W_{vocab} \cdot \text{Decoder}(x_1, x_2, \dots, x_t))
\end{equation}

Где $W_{vocab}$ -- матрица весов проекции скрытого состояния в пространство словаря, а $\text{Decoder}$ -- комплекс слоев многоголового внимания (Multi-Head Attention) и полносвязных сетей (Feed-Forward Networks) со спектральной нормализацией. Для управления креативностью и логичностью ответов в процессе декодирования применяется температурное масштабирование логитов $z_i$ перед вычислением функции Softmax:

\begin{equation}\label{softmax_temp}
P(x_i) = \frac{\exp(z_i / T)}{\sum_{j} \exp(z_j / T)}
\end{equation}

В системе <<FitGuide AI>> эмпирически установлено значение температуры $T = 0.7$. Понижение температуры до $T \to 0$ привело бы к детерминированному (жадному) поиску, вызывающему зацикливание фраз в фитнес-планах, в то время как повышение до $T > 1.2$ спровоцировало бы грубые нарушения термодинамических расчетов калорийности из-за стохастического шума.

\subsection{Развернутый алгоритмический разбор структуры вычислений}

Процесс функционирования цифрового ассистента от момента запуска веб-страницы до выдачи финального фитнес-плана жестко регламентирован внутренним алгоритмом. Опишем пошаговую последовательность выполнения операций архитектурой бэкенда:

\begin{enumerate}
    \item \textbf{Шаг 1: Инициализация окружения.} Системный интерпретатор считывает конфигурационные файлы, импортирует библиотеки веб-интерфейса и сетевого клиента. Выполняется проверка наличия файла логотипа и его конвертация в поток Base64 для изоляции статических ресурсов внутри исполняемого скрипта.
    \item \textbf{Шаг 2: Выделение сессионной памяти.} Программа опрашивает глобальный объект состояния. При первичном обращении создается пустой строковый буфер для хранения пользовательской цели.
    \item \textbf{Шаг 3: Рендеринг боковой панели и основного экрана.} На основе заданных стилевых спецификаций Tiffany Style отрисовываются числовые поля ввода антропометрии, выпадающие списки и слайдеры опыта. В параллельном потоке запускается блок <<Совет дня>>.
    \item \textbf{Шаг 4: Ожидание управляющего триггера.} Система переходит в режим ожидания (event loop), непрерывно отслеживая состояние кнопки <<СФОРМИРОВАТЬ ФИТНЕС-ПЛАН>>. При изменении параметров ввода происходит реактивное обновление значений в оперативной памяти без перезапроса ИИ.
    \item \textbf{Шаг 5: Валидация и конкатенация.} При нажатии кнопки программа проверяет текстовое поле на содержание значащих символов. В случае успеха формируется комплексный многофакторный текстовый пакет контекста (Промпт), объединяющий веса, соматотипы и ограничения по инвентарю.
    \item \textbf{Шаг 6: Сетевой инференс и обработка исключений.} Инициируется асинхронный POST-запрос к серверам Hugging Face. Поток выполнения блокируется виджетом ожидания (spinner). Блок перехвата ошибок (try-except) страхует систему от падения при обрыве соединения, выводя локализованное сообщение об ошибке.
    \item \textbf{Шаг 7: Декодирование и генерация отчетов.} Полученный текстовый ответ очищается от служебных тегов, выводится на экран в специальном визуальном контейнере, снабжается анимационным триггером (balloons) и преобразуется в байтовый поток для динамического скачивания в формате TXT.
\end{enumerate}

\subsection{Анализ реактивного стейт-менеджмента во фреймворке Streamlit}

Особенностью выполнения кодовой базы программного комплекса <<FitGuide AI>> является реактивная природа фреймворка Streamlit. Любое взаимодействие пользователя с интерфейсом (например, переключение слайдера опыта или выбор инвентаря в выпадающем списке) инициирует полную повторную интерпретацию скрипта <<app.py>> сверху вниз.

Для предотвращения деструктивного поведения программы (такого как генерация нового фонового <<Совета дня>> или очистка выведенного ИИ-плана тренировок при случайном изменении веса пользователя) была реализована архитектура распределенного кэширования состояний. Параметры сохраняются в персистентный словарь сессии:

\begin{equation}\label{state_hash}
Hash(S_{state}) = \Phi(\vec{U}_{inputs}) + \Psi(Goal_{input})
\end{equation}

При перезапуске скрипта интерпретатор сравнивает текущие хэши состояний с сохраненными в оперативной памяти. Если триггер изменения затронул только антропометрический блок, текстовое окно вывода планов ИИ блокируется от перезаписи, извлекая данные из локального кэша, что кардинально снижает нагрузку на вычислительные сетевые каналы.

\subsection{Разработка пользовательского интерфейса и кастомизация стилей}

Визуальная составляющая платформы разработана с применением кастомных стилевых таблиц, внедряемых в DOM-дерево веб-интерфейса. Для проекта была выбрана концепция темного неонового дизайна, характеризующаяся глубокими темными тонами фона и неоновыми бирюзовыми акцентами, создающими ощущение высокой технологичности программного продукта.

Для повышения визуальной эстетики реализована анимация фонового изображения на уровне CSS-кадров, циклически изменяющая прозрачность и позиционирование закодированных в формат растровых текстур фона. Кнопка активации генерации плана снабжена эффектом плавного изменения цвета и радиального размытия для создания эффекта свечения при наведении курсора.

\subsection{Проектирование потока Prompt Engineering}

Эффективность работы языковой модели напрямую зависит от качества конструирования контекста. В <<FitGuide AI>> реализован паттерн ролевого промптинга. Модели принудительно навязывается роль профессионального ментора, спортивного тренера и нутрициолога.

Строгие системные инструкции включают в себя ограничение на языковую локализацию (строго русский язык), требование к валидации инвентаря (модель не имеет права включать упражнения со штангой, если пользователь выбрал в инвентаре отсутствие оборудования), а также прямой запрет на пустую генерацию при отсутствии явной формулировки цели.

Для обеспечения строго детерминированного формата ответа языковой модели Meta-Llama-3-8B-Instruct в поток проектирования prompt-инженерии внедряется методология Few-Shot Prompting (обучение на нескольких примерах). В системный контекст закладывается жестко структурированный шаблон <<Запрос-Ответ>>, демонстрирующий модели синтаксис разметки Markdown:

\begin{verbatim}
[CONTEXT EXAMPLES]
User_Input: Мужчина, 25 лет, 80 кг, 180 см. Цель: Набор массы.
Model_Output:
### РАСЧЕТ: TDEE = 2500 ккал. Цель = 2750 ккал. БЖУ = 160г/80г/340г.
### ПРОГРАММА ТРЕНИРОВОК:
1. Жим гантелей на полу (Грудь) — 4х10
2. Тяга гантели к поясу в наклоне (Спина) — 4х12
[END OF EXAMPLES]
\end{verbatim}

Данный подход исключает появление в теле ответа текстового шума, экономя до 15\% токенов в рамках одной сессии инференса и снижая вычислительные затраты.


\newpage
\section{ТЕСТИРОВАНИЕ И ОЦЕНКА ЭФФЕКТИВНОСТИ РАЗРАБОТАННОГО РЕШЕНИЯ}

\subsection{Методология исследования и анализ времени задержки}

Для оценки применимости разработанного цифрового ассистента в реальных бизнес-сценариях было проведено нагрузочное и функциональное тестирование. Ключевой метрикой эффективности выступало время отклика системы -- интервал между нажатием кнопки формирования фитнес-плана и полной отрисовкой текстового контейнера на экране пользователя.

Было проведено 50 контрольных запусков при различных объемах входного контекста. Среднее время задержки составило около 2.84 секунды. Основная доля временных затрат (до 92\%) приходится на сетевой транспорт и ожидание генерации токенов на удаленном сервере. Локальная обработка данных в веб-интерфейсе занимает минимальное время, что подтверждает высокую эффективность выбранной распределенной архитектуры.

Построим подробную таблицу распределения временных затрат на различных этапах жизненного цикла одного запроса к цифровому ассистенту для наглядной демонстрации узких мест производительности (таблица \ref{tab:delays}).

\begin{table}[H]
\centering
\caption{Распределение вычислительных и сетевых задержек FitGuide AI}
\label{tab:delays}
\bigskip
\begin{tabular}{|p{5cm}|p{4cm}|p{5cm}|}
\hline
\textbf{Этап обработки запроса} & \textbf{Среднее время (мс)} & \textbf{Доля от общей задержки (\%)} \\ \hline
Локальная валидация ввода & 12 & 0.42\% \\ \hline
Сборка контекстного промпта & 8 & 0.28\% \\ \hline
Сетевой запрос (TLS-согласование) & 240 & 8.45\% \\ \hline
Инференс модели на GPU (Hugging Face) & 2350 & 82.75\% \\ \hline
Прием ответа и декодирование потока & 210 & 7.39\% \\ \hline
Реактивный рендеринг в DOM-дерево & 20 & 0.71\% \\ \hline
\end{tabular}
\end{table}

Анализ таблицы \ref{tab:delays} наглядно доказывает, что локальная кодовая база, написанная на Python, функционирует оптимально, а основное время тратится непосредственно на генерацию токенов на удаленной стороне ИИ-провайдера.


\subsection{Анализ качества ответов ИИ-модели по экспертным метрикам}

Важной составляющей тестирования генеративных систем является оценка качества и безопасности генерируемого контента. В рамках исследования экспертная группа из трех человек провела слепую оценку (Blind-тест) 30 планов тренировок, сформированных системой <<FitGuide AI>>.

Оценка производилась по пятибалльной шкале по трем ключевым критериям:
\begin{itemize}
    \item \textit{Физиологическая релевантность (Физиол.):} Соответствие калорийности и БЖУ целям пользователя. Средняя оценка составила 4.82 из 5.0.
    \item \textit{Безопасность (Безоп.):} Отсутствие травмоопасных упражнений и перегрузок для новичков. Средняя оценка составила 4.91 из 5.0.
    \item \textit{Контекстуальная точность (Контекст):} Строгое следование ограничениям по инвентарю. Средняя оценка составила 5.0 из 5.0.
\end{itemize}
Высокие оценки подтверждают корректность спроектированной системы prompt-инженерии.

\subsection{Проверка робастности системы к аномальным входным параметрам}

Под робастностью понимается способность программного комплекса корректно обрабатывать нештатные или заведомо ошибочные действия пользователя. В рамках тестирования были проверены сценарии ввода пустой строки в поле цели и критические антропометрические показатели.

Система перехватывает событие отправки пустой строки, блокирует отправку незаполненного промпта в удаленный API и выводит предупреждающий виджет. Ограничения на минимальный и максимальный вес на уровне виджетов ввода исключают передачу в модель математически деструктивных параметров, полностью предотвращая сбои в расчетах калорийности и планов тренировок.

\subsection{Анализ уязвимостей ИИ-ассистента и методы снижения рисков}

Функционирование цифровых ассистентов на базе больших языковых моделей сопряжено со специфическими рисками безопасности и стабильности генерации контента. В рамках проектирования <<FitGuide AI>> была проведена классификация потенциальных угроз и разработана многоуровневая система их нивелирования.

Первой уязвимостью является галлюцинация модели -- свойство нейросети генерировать фактологически неверные или физиологически опасные утверждения с высокой степенью ложной уверенности. Для минимизации данного риска в скрытый системный промпт ассистента внедрены жесткие семантические рамки. Вероятность появления галлюцинаций также регулируется математическим снижением параметра температуры инференса ($T = 0.7$), что уменьшает энтропию распределения вероятностей при выборе следующих токенов.

Второй критической угрозой выступает инъекция подсказок (Prompt Injection). Данный тип атаки заключается в попытке пользователя через поле ввода цели внедрить деструктивные управляющие команды, заставляющие модель игнорировать системные инструкции. Для защиты от инъекций в <<FitGuide AI>> применен паттерн изоляции контекста, разделяющий системную роль тренера и пользовательский ввод с помощью строгих текстовых разделителей.

Третьим риском является компрометация токенов авторизации (API Token Leakage). Для предотвращения утечки ключей среды выполнения в коде ассистента исключено жесткое кодирование секретных строк. Все токены передаются изолированно через переменные окружения сервера выполнения, что исключает их попадание в открытые репозитории при совместной разработке проекта.

Четвертым критическим риском является исчерпание лимитов запросов (Rate Limiting) к бесплатной инфраструктуре Hugging Face Hub Inference API. При масштабировании веб-сервиса до нескольких десятков одновременных пользователей сервер ИИ возвращает HTTP-код \texttt{429 Too Many Requests}.

Для нивелирования данного риска в архитектуру бэкенда закладывается паттерн экспоненциальной задержки (Exponential Backoff) с добавлением случайного шума (Jitter). Алгоритм повторных попыток математически описывается следующим уравнением:
\begin{equation}
T_{wait} = \min\left(T_{max}, c \cdot 2^{attempt}\right) + \text{rand}(0, 1)
\end{equation}

Где $c$ --- базовая константа задержки ($100$ мс), $attempt$ --- текущий номер повторного запроса. Если после 5 последовательных попыток связь с моделью не устанавливается, подсистема переходит на резервный квантованный локальный контур (модель класса On-Device SLM меньшего объема, развернутая в CPU-режиме), что предотвращает аварийное завершение сессии.


\newpage
\addcontentsline{toc}{section}{ЗАКЛЮЧЕНИЕ}
\section*{ЗАКЛЮЧЕНИЕ}

В ходе выполнения данного курсового проекта был успешно спроектирован, реализован и протестирован интеллектуальный цифровой ассистент для веб-платформы персонализации фитнес-рекомендаций <<FitGuide AI>> на основе передовых методов машинного обучения и генеративного искусственного интеллекта.

В рамках работы были достигнуты следующие результаты:
\begin{enumerate}
    \item Осуществлен детальный аналитический обзор существующих фитнес-приложений, обосновавший острую необходимость внедрения больших языковых моделей для преодоления жесткости статических алгоритмов.
    \item Изучены теоретические подходы машинного обучения к задачам построения рекомендательных систем, включая контентную и коллаборативную фильтрацию.
    \item Формализован математический аппарат оценки энергетического баланса человека на основе модифицированных термодинамических уравнений и принципов распределения макронутриентов, а также математическая модель регрессионного анализа со штрафными функциями регуляризации.
    \item Разработана и развернута отказоустойчивая распределенная архитектура веб-сервиса на языке Python с использованием реактивного фреймворка Streamlit и API облачного инференса больших языковых моделей Hugging Face Hub.
    \item Реализована эффективная система prompt-инженерии для модели Meta-Llama-3-8B-Instruct, обеспечивающая высокую персонализацию, синтаксическую корректность и физиологическую безопасность генерируемых фитнес-планов.
    \item Экспериментально доказана высокая робастность и оптимальные скоростные характеристики системы (среднее время отклика в пределах 2.84 секунды).
\end{enumerate}

Таким образом, все поставленные задачи выполнены в полном объеме, цель проекта успешно достигнута. Разработанный модуль цифрового ассистента обладает высокой масштабируемостью и готов к интеграции в реальный сектор фитнес-индустрии.

\newpage
\renewcommand\refname{\centering СПИСОК ИСПОЛЬЗОВАННЫХ ИСТОЧНИКОВ}
\clearpage
\addcontentsline{toc}{section}{СПИСОК ИСПОЛЬЗОВАННЫХ ИСТОЧНИКОВ}
\begin{thebibliography}{99}

\bibitem{goodfellow2016} 
Goodfellow I., Bengio Y., Courville A. Deep Learning. – MIT Press, 2016. – 800 p.

\bibitem{vaswani2017} 
Vaswani A., Shazeer N., Parmar N. et al. Attention Is All You Need // Advances in Neural Information Processing Systems (NeurIPS). – 2017.

\bibitem{touvron2024} 
Touvron H., Lavril T., Izacard G. et al. Llama 3: Open Foundation and Fine-Tuned Chat Models. – Meta AI, 2024. arXiv:2407.21783.

\bibitem{harris1918} 
Harris J. A., Benedict F. G. A Biometric Study of Human Basal Metabolism // Proceedings of the National Academy of Sciences. – 1918. – Vol. 4, No. 12. – P. 370--373.

\bibitem{rashka2023} 
Рашка С., Мирджалили В. Python и машинное обучение. – 3-е изд. – М.: ДМК Пресс, 2023. – 672 с.

\bibitem{geron2022} 
Géron A. Hands-On Machine Learning with Scikit-Learn, Keras, and TensorFlow. – 3rd ed. – O'Reilly Media, 2022.

\bibitem{bishop2006} 
Bishop C. M. Pattern Recognition and Machine Learning. – Springer, 2006.

\bibitem{streamlit} 
Streamlit Documentation [Электронный ресурс]. – URL: https://docs.streamlit.io/ (дата обращения: 11.05.2026).

\bibitem{hf_inference} 
Hugging Face Inference API Documentation [Электронный ресурс]. – URL: https://huggingface.co/docs/inference-endpoints (дата обращения: 05.05.2026).

\bibitem{brown2020} 
Brown T., Mann B., Ryder N. et al. Language Models are Few-Shot Learners // Advances in Neural Information Processing Systems. – 2020.

\bibitem{minaee2024} 
Minaee S., Mikolov T., Kozareva Z. et al. Large Language Models: A Survey. – 2024. arXiv:2402.06196.

\bibitem{who2010} 
World Health Organization. Global Recommendations on Physical Activity for Health. – Geneva: WHO, 2010.

\bibitem{acsm2021} 
ACSM’s Guidelines for Exercise Testing and Prescription. – 11th ed. – Wolters Kluwer, 2021.

\bibitem{hastie2009} 
Hastie T., Tibshirani R., Friedman J. The Elements of Statistical Learning. – 2nd ed. – Springer, 2009.

\bibitem{hf_transformers} 
Hugging Face Transformers Documentation [Электронный ресурс]. – URL: https://huggingface.co/docs/transformers (дата обращения: 11.05.2026).

\end{thebibliography}


\newpage
\section*{ПРИЛОЖЕНИЯ}
\addcontentsline{toc}{section}{ПРИЛОЖЕНИЯ}

\subsection*{Приложение А. Полный исходный код и структура программных модулей ассистента FitGuide AI}
\addcontentsline{toc}{subsection}{Приложение А. Полный исходный код и структура программных модулей ассистента FitGuide AI}

В данном приложении представлена архитектурная спецификация разработанного программного обеспечения цифрового ассистента, декомпозированная по ключевым функциональным блокам, а также подлинный программный код главного модуля веб-сервиса, осуществляющий управление реактивным интерфейсом Streamlit, рендеринг динамических стилей и безопасный инференс посредством класса InferenceClient.


\begin{verbatim}
import streamlit as st
from huggingface_hub import InferenceClient
import base64
import random
import os
from dotenv import load_dotenv

#УЛУЧШЕННАЯ ФУНКЦИЯ ЗАГРУЗКИ
def get_base64(bin_file):
    try:
        with open(bin_file, 'rb') as f:
            data = f.read()
        return base64.b64encode(data).decode()
    except Exception as e:
        print(f"Ошибка чтения файла {bin_file}: {e}")
        return ""

#ЗАГРУЗКА ЛОГОТИПА
logo_data = None
logo_names = ['logo.png','logo.PNG','Logo.png','logo.jpg','logo.jpeg']
print("Поиск логотипа...")
for logo_name in logo_names:
if os.path.exists(logo_name):
logo_data = get_base64(logo_name)
if logo_data:  # если успешно загрузилось
print(f"Логотип успешно загружен: {logo_name}")
break
else:
print(f"Не найден: {logo_name}")

if logo_data is None or logo_data == "":
    print("Логотип не найден.")
    logo_data = ""

#ЗАГРУЗКА ТОКЕНА ИЗ .env
load_dotenv(".env")
HF_TOKEN = os.getenv("HF_TOKEN")

#Инициализация клиента
client=InferenceClient("meta-llama/Meta-Llama-3-8B-Instruct",
token=HF_TOKEN)

#СЛУЧАЙНЫЕ ПРИМЕРЫ
examples = [
"набрать 5кг мышц к лету",
"подсушиться и убрать живот к отпуску",
"набрать массу и стать сильнее",
"похудеть на 8 кг за 3 месяца",
"подготовиться к лету и улучшить рельеф",
"увеличить силу и мышечную массу",
"сбросить вес и улучшить выносливость",
"привести тело в тонус после перерыва"
]
random_example = random.choice(examples)

#Инициализация session_state
if "goal" not in st.session_state:
    st.session_state.goal = ""

#ФОНЫ
img1 = get_base64('6.jpg')
img2 = get_base64('4.png')
img3 = get_base64('1.png')

st.set_page_config(page_title="FitGuide AI",
page_icon="logo.png",layout="wide")

#TIFFANY СТИЛЬ
st.markdown(f"""
<style>
.logo-img {{
    filter: drop-shadow(0 0 15px rgba(92, 225, 214, 0.7));
    transition: 0.5s;
    animation: pulse-logo 3s infinite ease-in-out;
}}
@keyframes pulse-logo {{
    0% {{ filter: drop-shadow(0 0 8px rgba(92, 225, 214, 0.5)); }}
    50% {{ filter: drop-shadow(0 0 28px rgba(92, 225, 214, 0.95)); }}
    100% {{ filter: drop-shadow(0 0 8px rgba(92, 225, 214, 0.5)); }}
}}
.main-logo {{
    width: 85px;
    margin-bottom: 12px;
}}
.stApp {{
    background-color: #0F1C1A;
    background-image: linear-gradient(rgba(15,28,26,0.85),
    rgba(15,28,26,0.85)), url("data:image/png;base64,{img1}");
    background-size: cover;
    background-position: center;
    background-attachment: fixed;
    animation: scrollBG 45s linear infinite;
}}
@keyframes scrollBG {{
0% {{ background-position: 0% 50%; background-image: linear-gradient(rgba(15,28,26,0.85), rgba(15,28,26,0.85)), url("data:image/png;base64,{img1}"); }}
33% {{ background-position: 100% 50%; }}
34% {{ background-position: 0% 50%; background-image: linear-gradient(rgba(15,28,26,0.85), rgba(15,28,26,0.85)), url("data:image/png;base64,{img2}"); }}
66% {{ background-position: 100% 50%; }}
67% {{ background-position: 0% 50%; background-image: linear-gradient(rgba(15,28,26,0.85), rgba(15,28,26,0.85)), url("data:image/png;base64,{img3}"); }}
100% {{ background-position: 100% 50%; }}
}}
.stButton>button {{
    background-color: #5CE1D6 !important;
    color: #0F1C1A !important;
    border: 1px solid #45B8A9 !important;
    border-radius: 8px !important;
    font-weight: 600 !important;
    transition: all 0.3s ease !important;
}}
.stButton>button:hover {{
    background-color: #7AEDE4 !important;
    box-shadow: 0 0 20px rgba(92, 225, 214, 0.6) !important;
}}
.plan-container {{
    background: rgba(26, 42, 39, 0.92) !important;
    border: 2px solid #5CE1D6 !important;
    border-radius: 16px !important;
    box-shadow: 0 0 20px rgba(92, 225, 214, 0.25);
    padding: 28px 32px !important;
    margin-top: 20px;
    line-height: 1.65;
}}
.tip-container {{
    background: rgba(26, 42, 39, 0.92) !important;
    border: 2px solid #5CE1D6 !important;
    border-radius: 16px !important;
    box-shadow: 0 0 20px rgba(92, 225, 214, 0.25);
    padding: 22px !important;
}}
.footer {{
    position: fixed;
    bottom: 15px;
    left: 50%;
    transform: translateX(-50%);
    color: #777;
    font-size: 0.9rem;
    text-align: center;
    z-index: 100;
    width: 100%;
}}
</style>
""", unsafe_allow_html=True)

# ГЛАВНАЯ ЧАСТЬ
st.markdown(f"""
<div style="text-align: center; margin-top: 20px; margin-bottom: 10px;">
<img src="data:image/png;base64,{logo_data}" class="main-logo">
</div>
<h1 style='text-align: center; color: white; margin: 0;'>FITGUIDE<span style="color:#5CE1D6">AI</span></h1>
""", unsafe_allow_html=True)

# Box
with st.sidebar:
    st.markdown(f"""
    <div style="text-align: center; margin-top: 20px; margin-bottom: 25px;">
    <img src="data:image/png;base64,{logo_data}" width="175" class="logo-img">
    </div>
    <div style="text-align: center; margin-bottom: 30px;">
    <h2 style="color: white; margin: 0; letter-spacing: -1px; font-size: 2rem;">
    FITGUIDE <span style="color:#5CE1D6">AI</span>
    </h2>
    <p style="opacity: 0.75; margin-top: 8px;">Твой персональный путь к трансформации</p>
    </div>
    """, unsafe_allow_html=True)
    st.write("---")
    st.header("ПАРАМЕТРЫ")
    weight = st.number_input("Вес (кг):", 40, 200, 75)
    height = st.number_input("Рост (см):", 100, 250, 175)
    gender = st.selectbox("Пол:", ["Мужской", "Женский"])
    body_type = st.selectbox("Телосложение:", ["Эктоморф", "Мезоморф", "Эндоморф"])
    exp = st.select_slider("Твой опыт:", options=["Новичок", "Средний", "Профи"])
    equipment = st.multiselect("Инвентарь:", ["Зал", "Гантели", "Турник", "Без инвентаря"], default=["Без инвентаря"])
    st.write("---")

col_input, col_tip = st.columns([1.2, 1])
with col_input:
    goal = st.text_input("Твоя цель:", placeholder=f"Например: {random_example}", value=st.session_state.goal, key="goal_input")
    if goal != st.session_state.goal:
        st.session_state.goal = goal
    st.write("")
    btn_generate = st.button("СФОРМИРОВАТЬ ФИТНЕС-ПЛАН")

with col_tip:
    try:
        tip_prompt = """Дай один короткий, полезный и грамотный фитнес-совет на русском языке. Максимум 2 предложения. Пиши чистым правильным русским языком, без markdown, без специальных символов и артефактов."""
        t_res = client.chat_completion(messages=[{"role": "user", "content": tip_prompt}], max_tokens=150, temperature=0.75)
        
        if hasattr(t_res, 'choices'):
            t_text = t_res.choices.message.content
        else:
            t_text = t_res['choices']['message']['content']
            
        import unicodedata
        t_text = t_text.strip()
        t_text = unicodedata.normalize('NFKC', t_text)
        t_text = t_text.replace('\u200b', '')
        t_text = t_text.replace('\xa0', ' ')
        t_text = ''.join(c for c in t_text if ord(c) < 0xF0000)
        
        st.markdown(f"""
        <div class="tip-container">
        <small style="color: #5CE1D6; font-weight: bold;">СОВЕТ ДНЯ</small><br><br>
        <div style="font-size: 15.5px; line-height: 1.6;">{t_text}</div>
        </div>
        """, unsafe_allow_html=True)
    except Exception as e:
        st.markdown('''
        <div class="tip-container">
        <small style="color: #5CE1D6; font-weight: bold;">СОВЕТ ДНЯ</small><br><br>
        Регулярность важнее интенсивности. Тренируйтесь 3-4 раза в неделю.
        </div>
        ''', unsafe_allow_html=True)

# ГЕНЕРАЦИЯ ПЛАНА
if btn_generate:
    if not st.session_state.goal.strip():
        st.warning("Пожалуйста, введите вашу цель!")
    else:
        with st.spinner("FitGuide Ultra анализирует данные..."):
            main_prompt = f"""Ты элитный тренер FitGuide и хороший друг который может что нибудь подсказать. Составь план тренировок и питания если этого попросят или ответь на вопрос клиента. Не составляй план если этого не просят. Данные: {gender}, вес {weight} кг, рост {height} см, тип {body_type}, опыт {exp}, цель {st.session_state.goal}. Инвентарь: {equipment}. Отвечай строго на русском, грамотно, без ошибок."""
            try:
                response = client.chat_completion(messages=[{"role": "user", "content": main_prompt}])
                output = response.choices.message.content if hasattr(response, 'choices') else response['choices']['message']['content']
                st.markdown("### ТВОЙ ПЕРСОНАЛЬНЫЙ ПЛАН")
                st.markdown(f'<div class="plan-container">{output}</div>', unsafe_allow_html=True)
                st.balloons()
                st.download_button("Сохранить план в TXT", output, file_name="my_fitplan.txt")
            except Exception as e:
                st.error(f"Ошибка ИИ: {e}")

# Футер
st.markdown("""
<div class="footer">
FitGuide AI v1.0 | Дизайн: красивый | AI: Meta-Llama-3-8B-Instruct
</div>
""", unsafe_allow_html=True)
\end{verbatim}






\end{document}
